\section{Formalization}\label{app:formal}

\subsection{Statements and scope}
The development uses Cubical Agda \cite{cubicalagda} and the Cubical library \cite{cubical}. The entry \path{agda/src/DeloopS3/Assembly/Complete.agda} exports
\[
\texttt{Deloop-S3}:
\sum_{B:\mathcal U_{0,*}}
  \bigl(\Conn(B)\times(\Omega B\eqp S^3)\bigr).
\]
Here $\mathcal U_{0,*}$ is \texttt{Pointed} at level zero;
$\Conn(B)$ is \texttt{isConnected 2 (typ B)}, and the pointed sphere
is written $(\texttt{S}_{+}\;3,\mathsf{north})$ in the source.
The sum is untruncated. The term has no module parameters and uses no
postulates. The second entry, \path{agda/src/DeloopS3/HPInf/Complete.agda},
exports the pointed map \texttt{t}, its fibre equivalence \texttt{t-fiber},
and the combined declaration
\texttt{S\textsuperscript{3}-delooping-with-maximal-torus}.
The latter uses the same concrete $B$: its map is $K(\Z,2)\to_*B$
and its fibre over the basepoint is equivalent to $S^2$. It does not assert
an equivalence with an independently defined $\mathbb H P^\infty$.
The pointed loop equivalence is obtained by translating the source loop
space before applying the unpointed equivalence
(\texttt{point-loop-equivalence} in \path{Assembly/DefB.agda}); it does
not assert that the unmodified comparison preserves the basepoint or that
it respects the Cayley--Dickson multiplication.
Both entries use \texttt{--safe --cubical --guardedness}.
Some auxiliary modules use \texttt{--lossy-unification}, an elaboration
option with no additional axioms.

The formalized fracture theorem assumes $2$-connectedness, which suffices
for $F\simeq S^3$; the simply connected theorem cited in the paper is the
more general result of \cite[Thm.~5.4]{scoccola}. Moore recognition is
formalized for arbitrary positive $m$ and $2$-connected finite CW types
of dimension at most $4$ with the stated homology profile. The bottom-Moore
cofibre argument and the resulting contractions are formalized for the
$O$ and $I$ profiles used here. Classical recognition uses
\cite[Rem.~1.2]{dmw}, without a separate genus-membership assumption;
it is cited, not formalized. Independent $\mathbb H P^\infty$ and Cayley--Dickson
multiplication comparisons remain outside the exports.

The geometric Mayer--Vietoris interface supplies exactness for $0\leq n<5$,
with boundary $H_{n+1}\to H_n$. Its naturality comparisons establish
$H_3\to H_2$ naturality for maps between spans and $H_2\to H_1$ naturality
for self-span maps, retaining their comparison homotopies. These are the
degrees used in \S\ref{sec:homology}. The all-degree chain equivalence
and these degree-specific geometric comparisons have separate exports.

\medskip
\noindent\textbf{Artifact.} The formalization and accompanying paper are hosted at \url{https://github.com/wxkj99/s3-delooping}. The formalization source snapshot used for this paper is commit \href{https://github.com/wxkj99/s3-delooping/tree/2eda93321fa89daa4875e6921db3b5fae19f75f2}{\texttt{2eda933}}. The repository contains the formalization, checking tools, pinned Cubical dependency, and the paper's \LaTeX{} sources and PDF. Local audit documents and verification records are maintained separately.

\subsection{Indexing conventions}
The paper uses classical truncation degrees. The library uses the following
indices.
\begin{center}
\small
\begin{tabular}{@{}ll@{}}
\toprule
Mathematical notation & Library notation \\
\midrule
$X$ connected & \texttt{isConnected 2 X} \\
$X$ simply connected & \texttt{isConnected 3 X} \\
$X$ $2$-connected & \texttt{isConnected 4 X} \\
$\tau_1X$ & \texttt{hLevelTrunc 3 X} \\
$\tau_2X$ & \texttt{hLevelTrunc 4 X} \\
\bottomrule
\end{tabular}
\end{center}
The source notation $\texttt{PT.}\lVert X\rVert_1$ denotes propositional
truncation. The declaration \texttt{Moore m k} counts suspensions from
$S^1$: it denotes $\Moore{m}{k+1}$ in this paper. In particular,
\texttt{Moore 2 1} and \texttt{Moore 10 2} denote $\Moore22$ and
$\Moore{10}3$.

\subsection{Statement-to-declaration correspondence}
Module paths below are relative to \path{agda/src/DeloopS3/}, except that
paths starting with \path{Cubical/} are relative to
\path{vendor/cubical/}. Declaration names include any enclosing named
parameterized module. Each listed file is in the active source tree.
\begingroup
\small
\setlength{\tabcolsep}{4pt}
\begin{longtable}{@{}>{\raggedright\arraybackslash}p{0.29\textwidth}>{\raggedright\arraybackslash}p{\dimexpr0.71\textwidth-8pt\relax}@{}}
\toprule
Paper statement & Exported declaration and module \\
\midrule
\endfirsthead
\toprule
Paper statement & Exported declaration and module (continued) \\
\midrule
\endhead
\bottomrule
\endfoot
Theorem~\ref{thm:main} & \texttt{Deloop-S3}\newline
  \path{Assembly/Complete.agda} \\[3pt]
Truncation square & \texttt{TruncationSquare.cartesian}\newline
  \path{Sphere/TruncationCartesian.agda} \\[3pt]
Hopf reflection & \texttt{hopfReflection}, \texttt{sphereUnitFiber}\newline
  \path{Sphere/HopfUnit.agda} \\[3pt]
Common fibre $F\simeq S^3$ & \texttt{RP\textsuperscript{2}-unit-fiber}\newline
  \path{Sphere/RP2UnitFiber.agda} \\[3pt]
Finite covers and total-space equivalences & \texttt{short-total}, \texttt{p}, \texttt{p-fibres}, \texttt{p-is-cover}\newline
  \path{Cert/Interpret/Done.agda} \\[3pt]
Twisted circle bundle & \texttt{Bundle.circleFibres}, \texttt{Bundle.baseFibre}\newline
  \path{Seifert/TwistedBundle/Cell.agda} \\[3pt]
Clutching coefficient $-2e$ in the calibrated chart & \texttt{Calibrated.coefficient}, \texttt{Calibrated.calibratedTotal}\newline
  \path{Seifert/TwistedBundle/TwoCap.agda} \\[3pt]
$\pi_2(M_e)=0$ for $e\ne0$ & \texttt{pi\textsubscript{2}-M\textsubscript{e}-zero}\newline
  \path{Seifert/TwistedBundle/Pi2Zero.agda} \\[3pt]
Dihedral-character rigidity & \texttt{dihedral-character-rigid}\newline
  \path{Bridge/SpanPi1/Character.agda} \\[3pt]
Source character & \texttt{source-character}\newline
  \path{Bridge/SpanPi1/SourceCharacter.agda} \\[3pt]
Nonabelian image and kernel witness & \texttt{subgroup-image-nonabelian}, \texttt{normal-core-lift}\newline
  \path{Bridge/SpanPi1/CoverWitnesses.agda} \\[3pt]
Identification of the kernel character & \texttt{cover-kernel-character-from-source}\newline
  \path{Bridge/SpanPi1/KernelCharacter.agda} \\[3pt]
Simple connectedness of the total pushouts & \texttt{X-simple}\newline
  \path{Bridge/SpanPi1/Target.agda} \\[3pt]
Cartesian span squares & \texttt{left-cartesian}, \texttt{right-cartesian}\newline
  \path{Assembly/MixedTruncation.agda} \\[3pt]
Fibre sequences $F\to E_c\to P_c$ & \texttt{local-fibre}\newline
  \path{Assembly/LocalCubes.agda} \\[3pt]
Actual cellular homology & \texttt{CW-AugChainComplex}, $\widetilde{\texttt H}^{\texttt{skel}}$\newline
  \path{Cubical/CW/ChainComplex.agda};\newline
  $\widetilde{\texttt H}^{\texttt{cw}}$\newline
  \path{Cubical/CW/Homology/Base.agda} \\[3pt]
Actual Hurewicz map and theorem & \texttt{HurewiczMapUntrunc}, \texttt{HurewiczTheorem}\newline
  \path{Cubical/CW/HurewiczTheorem.agda} \\[3pt]
Geometric exactness, $n<5$ & \texttt{SpanSequenceData}, \texttt{GeometricPushoutLES}\newline
  \path{Bridge/DMCExactness/Geometric.agda};\newline
  \texttt{finite-pushoutLES-merely}\newline
  \path{Bridge/CellularPushout/PrismChains/Packaging/Unconditional.agda} \\[3pt]
Boundary naturality, $H_3\to H_2$ & \texttt{TotalWindowBoundary.geometric-top-boundary-natural}\newline
  \path{Bridge/DegreeSpan/ContinuousBoundaryNatural.agda};\newline
  \texttt{initial-boundary-natural}\newline
  \path{Bridge/DegreeSpan/InitialBoundaryNatural.agda} \\[3pt]
Self-span naturality, $H_2\to H_1$ & \texttt{finite-same-span-boundary-natural}\newline
  \path{Bridge/DegreeSpan/LowerBoundaryDescent.agda} \\[3pt]
Component homology & \texttt{component-H\textsubscript{1}-groups},
  \texttt{source-H\textsubscript{2}-zero}, \texttt{target-H\textsubscript{2}-zero},
  \texttt{source-H\textsubscript{3}}, \texttt{target-H\textsubscript{3}}\newline
  \path{Bridge/Profile/ComponentHomologyComplete.agda} \\[3pt]
Covering-degree comparison & \texttt{comparisons}\newline
  \path{Bridge/DegreeSpan/LeftComparisons.agda};\newline
  \texttt{cover-degree-from-H2}\newline
  \path{Bridge/DegreeSpan/CoverDegreeFromH2.agda} \\[3pt]
Four-dimensional Moore recognition & \texttt{recognize-moore3}, \texttt{moore3-map-isEquiv}\newline
  \path{Bridge/Profile/MooreRecognition.agda} \\[3pt]
Bottom cofibre for the $I$ profile & \texttt{bottom-cofib-profile}\newline
  \path{Bridge/Profile/BottomCofibre.agda} \\[3pt]
Local contractions & \texttt{O-profile-contraction}, \texttt{I-profile-contraction}\newline
  \path{Bridge/Profile/Contraction.agda} \\[3pt]
Fibre preservation over simply connected bases & \texttt{localization-preserves-fiber}, \texttt{fiber-comparison}\newline
  \path{Local/RouteB2.agda} \\[3pt]
Fracture for $2$-connected types & \texttt{arithmetic-fracture-2-connected}\newline
  \path{Local/Fracture/Fracture.agda} \\[3pt]
Assembly of the delooping & \texttt{B-def}\newline
  \path{Assembly/DefB.agda};\newline
  \texttt{Deloop-final}\newline
  \path{Assembly/Final.agda} \\[3pt]
Recognition of the specified $3/10$ fracture cone & \texttt{unit-aware-fracture}\newline
  \path{HPInf/UnitFracture.agda} \\[3pt]
Cartesianness of \eqref{diag:HP} & \texttt{common-square-cartesian}\newline
  \path{HPInf/RelativeFracture.agda} \\[3pt]
Maximal torus (Theorem~\ref{thm:torus}) & \texttt{t}, \texttt{t-fiber}, \texttt{maximal-torus}\newline
  \path{HPInf/Complete.agda} \\
\end{longtable}
\endgroup

\subsection{Reproduction and development size}
\paragraph{Dependency identity.}
The official Cubical commit is\newline
\path{87af036ebd1c87956cf08fea44e62c7cf76392ca}, with parent\newline
\path{3d828afcac588414920565ad22d08d3dcdf02249}.
The full unmodified snapshot, including original licenses, is in
\path{vendor/cubical/}, locked by \path{agda/deps.lock.json}; official
provenance was verified through the GitHub API. The identifier
\texttt{cubical-0.9} does not mean release \texttt{v0.9}, which lacks the
required CW homology and Hurewicz modules. \path{agda/src/DeloopS3/Compat/Whitehead.agda}
reexports three declarations from \path{Cubical.Homotopy.Whitehead},
without copied mathematical types or proofs. Five files have import-only changes;
proof bodies and vendor sources are unchanged. The closure has 467 project and 416 official modules (883 total, excluding compiler builtins).

\paragraph{Reproduction.}
The checking scripts require Agda 2.8.0.2, Python 3.9 or later, Bash,
\texttt{flock}, and GNU \texttt{timeout}. From the root of a repository checkout, run:
\begin{quote}\small
\begin{verbatim}
bash agda/check.sh --configure-only
DELOOP_TIMEOUT=18000 bash agda/check.sh
\end{verbatim}
\end{quote}
The first command verifies the compiler version and dependency manifest
and creates a machine-local library registry. The second checks both
entries listed in \path{agda/roots.txt}. The runner uses this registry
and disables default libraries. It resolves paths relative to the
checkout, supports paths containing spaces, and performs no downloads.
Each check runs with one RTS capability and a 12\,GiB heap limit;
\texttt{DELOOP\_TIMEOUT} sets the per-entry time limit in seconds.

\paragraph{Verification.}
The official-dependency build began in a relocated directory containing
spaces with an empty user home and no project or library interfaces.
Five focused checks and both final entries exited with code zero;
successive stages reused only interfaces generated within that run.
The complete sequence took about 56 minutes, with 3,109 seconds for the
assembly stage and 54 seconds for the maximal-torus stage. Peak recorded
child-process RSS was approximately 12.4\,GiB under a 12\,GiB GHC heap cap.
The configuration and audit tools pass 48 regression tests.
The repository sources preserve all non-comment program lines of the checked
source; their only differences are comment and blank-line cleanup.
Both cleaned entry points subsequently passed an incremental
recheck with exit code zero.
Commands, hashes and measurements are retained in the author's local
verification record, \path{docs/verification/verification.json}; this file
is not included in the source repository.

\paragraph{Semantic audit.}
The author's local audit record, \path{docs/formalization-audit.md},
records the checked statement-to-source interfaces, assumptions, and coverage
limits separately from compilation and portability status. No target substitution was found in those
interfaces. This is not an exhaustive independent review of the deep
cubical coherence proofs.

\paragraph{Size.}
The active tree \path{agda/src} contains 467 Agda modules and 64,352 lines,
including comments and blank lines; \path{agda/archive} is excluded.
The two generated certificate-data modules
\path{Cert/Octa3.agda} and \path{Cert/Icosa10.agda} account for 453 of
these lines (27,105 bytes). The modules under \path{Cert/Input/} implement
the geometric interpretation of the input data and are included in the
remaining proof source. The totals can be reproduced from the repository
root by the following one-line command, whose output is the module count
followed by the line count.
\begin{quote}\scriptsize
\begin{verbatim}
find agda/src -name '*.agda' -exec wc -l {} + | awk '$2!="total"{n++;s+=$1}END{print n,s}'
\end{verbatim}
\end{quote}

\subsection{Finite certificates}
The certificate data are included as Agda declarations in
\path{agda/src/DeloopS3/Cert/Octa3.agda} and
\path{agda/src/DeloopS3/Cert/Icosa10.agda}. They specify permutation tables,
coset representatives, Schreier spanning trees, boundary words, filling
multiplicities, and the cancellation and shear data used in the short covers.
Declarations such as \texttt{group-ok}, \texttt{cosets-ok},
\texttt{graph-ok}, and \texttt{normalization-ok} validate the finite
computations by reduction against \path{Cert/Check.agda}.

The modules under \path{Cert/Input/} connect these checked data to the
covering diagrams. Those under \path{Cert/Interpret/} construct equivalences
of higher inductive types and the required boundary comparisons.
\path{Cert/Interpret/Done.agda} exports \texttt{short-total}, \texttt{p},
and \texttt{p-fibres}. Agda checks both the finite certificate assertions
and this semantic interpretation used by the final theorem.
